\documentclass[10pt,a4paper]{amsart}
\usepackage[margin=19mm]{geometry}
\usepackage{amsmath,amssymb,mathtools}
\usepackage[scr=rsfs,cal=euler]{mathalfa}
\usepackage{xfrac}
\usepackage[unicode,hidelinks,bookmarks=false]{hyperref}
\newcommand{\ie}{i.e.\ }
\newcommand{\cf}{cf.\ }
\newcommand{\resp}{respectively\ }
\newcommand{\Subsection}{\subsection}
\providecommand{\nobreakdash}{\nobreak\hbox{-}}
\setlength{\emergencystretch}{2em}
\multlinegap0pt
\usepackage{bm}
\usepackage{bbm}
\usepackage{dsfont}
\usepackage{xspace}
\usepackage{cancel}
\usepackage{mathtools}
\usepackage{amsthm}
\makeatletter
\@ifundefined{claimproof}{\newtheorem{claimproof}{Claimproof}}{}
\@ifundefined{corollary}{\newtheorem{corollary}{Corollary}}{}
\@ifundefined{mclaim}{\newtheorem{mclaim}{Claim}}{}
\@ifundefined{theorem}{\newtheorem{theorem}{Theorem}}{}
\makeatother
\makeatletter
\expandafter\ifx\csname claimproof\endcsname\relax
\newenvironment{claimproof}
{\noindent {\em Proof of Claim:} }
{\hfill $\diamond$ \smallskip}
\fi
\newcommand{\forbiddenGraph}{F}
\newcommand{\sharedVertex}[1]{C^{\mathrm{V}}_{#1}}
\makeatother
\title[Bounding Width on Graph Classes of Constant Diameter]{Bounding Width on Graph Classes of Constant Diameter}
\author{Konrad K. Dabrowski, Tala Eagling-Vose, Noleen K\"ohler, Sebastian Ordyniak, Dani\"el Paulusma}
\date{Source: arXiv:2505.19926v1 (CC BY 4.0)}
\begin{document}
\maketitle

\noindent\emph{Source and licence.} Bounding Width on Graph Classes of Constant Diameter. Version arXiv:2505.19926v1 (CC BY 4.0), distributed under the Creative Commons Attribution 4.0 International licence, as recorded in the arXiv licence metadata for this exact version and shown on the abstract page https://arxiv.org/abs/2505.19926v1. Full rights notice: licenses/arxiv-2505.19926-CC-BY-4.0.txt.\par\smallskip

\noindent\textit{Editorial scope.} Counted unit: the Theorem whose source label is \texttt{thm:diam2-CVktimesl} in the article (source0.tex lines 1498--1501, source label \texttt{thm:diam2-CVktimesl}), delivered together with the article's own complete original proof of it (source0.tex lines 1502--1566, one \texttt{proof} environment of 65 lines). No mathematics is written, repaired or re-derived: statement and proof are the article's own text line for line. Required context retained uncounted: cor:longInducedPath (lines 327--330). Every definition, display and label of those lines is retained unchanged. Omitted from this package and not redistributed: 1--326, 331--1497, 1567--1728. Nothing inside a retained excerpt is omitted or altered: every retained body is delivered exactly as the article's lines are written, so no omission receipt of any kind accompanies this package. Citations: the retained excerpts carry no citation command at all, so no bibliography entry of the article is transcribed and no reference is redistributed. Reference adaptation: none was needed and none was made; every reference inside the delivered unit resolves inside it, so no unresolved reference or citation is produced. Printed slips kept literal: no printed word, symbol, hypothesis or number of the counted statement or of its proof was altered. Layout and class: the article's own class is \texttt{llncs} with packages including fontenc, graphicx, enumitem, amsmath, amssymb, ntheorem, todonotes, float, tkz-graph, hyperref, cleveref, url, booktabs; the wrapper is the lane's pinned \texttt{amsart} editorial wrapper at 10pt on A4, which restates the theorem-like environments the retained text needs and adds the article's own macro definitions that the retained text needs (\texttt{\textbackslash forbiddenGraph}, \texttt{\textbackslash sharedVertex}), with extra wrapper packages (bm, bbm, dsfont, xspace, cancel, mathtools). The delivered numbering is the unit's own. Assets: no retained span contains a figure environment, a graphics-inclusion command, a PDF-page inclusion, a table or a TikZ picture; no image file is copied into this package, the package declares no asset dependency and no image file is redistributed with it. Licence: version arXiv:2505.19926v1 of this article is distributed under the Creative Commons Attribution 4.0 International licence, as recorded in the arXiv licence metadata for this exact version and shown on the abstract page https://arxiv.org/abs/2505.19926v1. A full rights notice is delivered at licenses/arxiv-2505.19926-CC-BY-4.0.txt. Characters: every retained excerpt is pure ASCII, so the delivered engine, XeLaTeX with its default Latin Modern text font, reports no missing character.
\begin{corollary}\label{cor:longInducedPath}
    For all $r,s,\ell\in \mathbb{N}$, there is a number $c(r,s,\ell)$
    such that every $K_{r,s}$-subgraph-free graph of treedepth $c(r,s,\ell)$ has an induced $P_\ell$.
\end{corollary}

\begin{theorem}  
\label{thm:diam2-CVktimesl}
    For any integers $\ell\geq 3$ and $k\geq 1$ the class of all $\sharedVertex{k*[2\ell]}$-subgraph-free graphs of diameter at most~$2$ has bounded treedepth.
\end{theorem}

\begin{proof}
    Let $G$ be a $\sharedVertex{k*[2\ell]}$-subgraph-free graph of diameter at most~$2$, we claim $td(G) \leq c(k (2\ell-1)+1, k (2\ell-1),4^{k+1}k^{k+4}\ell)$. Suppose for contradiction $td(G) > c(k (2\ell-1)+1, k (2\ell-1),4^{k+1}k^{k+4}\ell)$. $G$ cannot contain a large complete bipartite subgraph as $K_{k (2\ell-1)+1, k (2\ell-1)}$ contains $\sharedVertex{k*[2\ell]}$ as a subgraph. From Corollary~\ref{cor:longInducedPath},  $G$ contains an induced path, $P=(p_0,\dots,p_m)$ 
    of length $m=4^{k+1}k^{k+4}\ell$. 

    We define the distance of two pairs  $(p_i,p_j)$ and $(p_{i'},p_{j'})$ of vertices on $P$ with $i<j$ and $i'<j'$ as follows. If either $i\leq i'\leq j$ or $i\leq j'\leq j$ we set the distance of $(p_i,p_j)$ and $(p_{i'},p_{j'})$ to be $0$. Otherwise, the distance of $(p_i,p_j)$ and $(p_{i'},p_{j'})$ is the positive integer $d$ for which $j'+d=i$ if $j'<i$ or $j+d=i'$ if $j<i'$.
    
    Suppose some vertex $v \in V(G)$ has $k$  pairs of neighbours in $P$, $(p_i,p_j)$ such that $j = i+2\ell-2$ and these $k$ pairs have pairwise distance at least $1$. Such a vertex results in \forbiddenGraph\ as there are $k$ cycles of length $2\ell$ each containing the single common vertex $v$. Hence, no vertex is adjacent to $k$ pairs of neighbours $(p_i,p_{i+2\ell-2})$ of pairwise distance at least $1$.

    \begin{mclaim}\label{clm:dist2l-4}
            Let $x$ be some vertex in $G$, then $N(x) \cap \{p_0,\dots,p_{m-\ell}\}$ contains at most $(k-1)^2$ disjoint pairs $\{p_i,p_j\}$ where $j = i+2\ell-4$, $j \leq m- \ell$ and each pair has pairwise distance at least $\ell$.
    \end{mclaim}
    \begin{claimproof}
            Say $x$ has neighbours $p_i, p_{i+2\ell-4} \in \{p_0,\dots,p_{m-\ell}\}$. The vertices $p_{i+\ell-3}$, $p_{i+3\ell-5}$ must have a common neighbour, $x'$. We call the vertex $x'$, that is the common neighbour of $p_{i+\ell-3}$ and $p_{i+3\ell-5}$ the \emph{connector} of the pair $(p_i,p_j)$ of neighbours of $x$. If $x'=x$ then there is some $C_{2\ell}$ containing only $x$ and the path vertices $p_{i+\ell-3}, \dots, p_{i+3\ell-5}$, otherwise $x'\neq x$ and there is a $C_{2\ell}$ given by $(x, p_i, \ldots, p_{i+\ell-3}, x', p_{i+3\ell-5}, \ldots, p_{i+2\ell-4})$. Note there is a cycle of length $2\ell$ containing $x'$ and the path vertices $p_{i+\ell-3}$, \ldots, $p_{i+3\ell-5}$.

            Suppose $N(x) \cap \{p_0,\dots,p_{m-\ell}\}$ contains $(k-1)^2+1$ disjoint pairs with pairwise distance at least $\ell$. Let $x'_r$ denote the \emph{connector} for the $r$th pair with $X' = \{x'_1,\dots, x'_{(k-1)^2+1}\}$. If $X'$ contains $k$ pairwise distinct vertices, then $G$ contains \forbiddenGraph\ with $k$ cycles of length $2\ell$ each cycle, apart from possibly one, containing $x$ and a distinct vertex $x' \in X'$. Note that one of these distinct connectors could be the vertex $x$ itself, in which case the cycle does not contain an additional vertex $x'\in X'$. 
            On the other hand, consider $|X'| \leq k-1$. As there are $(k-1)^k+1$ connector there must be some $x'\in X'$ which is the connector of at least $k$ distinct pairs. Assume without loss of generality, $x'_1=x'_2=\ldots=x'_k$. However, this is a contradiction as there are $k$ cycles of length $2\ell$ containing vertices from $P$ and the single common vertex $x'_1$.
    \end{claimproof}

    \begin{mclaim}\label{clm:vbndnei}
            For every $x \in V(G)$, there are less than $4^kk^{k+1}$ vertices in $N(x) \cap \{p_0,\dots,p_{m-3\ell-4}\}$ 
            with pairwise distance at least $3\ell-4$.
    \end{mclaim}
    \begin{claimproof}
        Consider some vertex $x_0$ and let $Z_0 \subseteq N(x_0) \cap \{p_0,\dots,p_{m-3\ell-4}\}$ be some set of vertices with pairwise distance at least $3\ell-4$ along the path. Suppose $|Z_0| \geq 4^kk^{k+1}$. 
        Let $Z^+_0 = \{ p_{i+2\ell-4}: p_i \in Z_0 \}$, note $|Z^+_0| = |Z_0|$.
        
         In the following we recursively construct vertices $x_0, x_1,\ldots,x_{k-1}$, sets $Z_{0} \supseteq Z_1 \supseteq \ldots \supseteq Z_{k-1}$ and sets $\hat{Z}_{0} \supseteq \hat{Z}_{1} \supseteq \ldots \supseteq \hat{Z}_{\delta-1} \supseteq \hat{Z}_{k-1}$ such that for every $1\leq i \leq k-1$
         \begin{enumerate}
             \item $|Z_i|\geq \left\lfloor {\frac{|Z_{i-1}|- (k-1)^2|}{2(k-1)}} \right\rfloor \geq \left\lfloor {\frac{|Z_{i-1}|}{4(k-1)}} \right\rfloor$, where $|Z_{i-1}| \geq 2(k-1)^2$;
             \item $|\hat{Z}_i|\geq |Z_i|-(k-1)^2$;
             \item $Z_i\subseteq N(x_0, \dots, x_i)\cap \{p_0,\dots, p_m\}$;
             \item $\hat{Z}_i\cap N(x_0,\dots, x_{i})=\emptyset$;
             \item $N(x_i)\cap \hat{Z}_{i-1}\neq \emptyset$; and 
             \item $\hat{Z}_i\subseteq Z_0^+$.
         \end{enumerate} 
        
        Suppose for $\delta \leq k-2$ we have constructed vertices $x_0,x_1,\ldots,x_{\delta}$ and sets $Z_{0} \supseteq Z_1 \supseteq  \ldots \supseteq Z_{\delta}$ and $\hat{Z}_0 \supseteq \hat{Z}_1 \supseteq \ldots \supseteq \hat{Z}_{\delta-1}$ with the properties above. Let $Z^+_{\delta} = \{ p_{i+2\ell-4}: p_i \in Z_{\delta-1}\}$. Again from Claim~\ref{clm:dist2l-4}, at most $(k-1)^2$ vertices in $Z^+_{\delta}$ are adjacent to $x_{\delta}$. Let $\hat{Z}_{\delta} =Z^+_{\delta} \setminus N(x,x_1, \ldots, x_{\delta})$ with $|\hat{Z}_{\delta}| \geq |Z_{\delta}| -(k-1)^2$. Thus properties $2$, $4$ and $6$ hold.
        Let $(p_i,p_j)$ be some pair such that $p_i \in Z_{\delta}$, $j\not=i+2\ell-4$ and $p_j \in \hat{Z}_{\delta}$. As $p_i,p_j$ must have some common neighbour, $x' \neq x$, there is some $C_{2\ell}$ given by $(x,p_i, x', p_j, \cdots, p_{j-(2\ell-4)})$. Note that we can choose  $\left\lfloor {\frac{|Z_{\delta-1}- (k-1)^2|}{2}} \right\rfloor$ such pairs with distance at least $1$.
        Let $(p_{r_1},p_{r_2})$ denote the $r$th pair and let $x'_r$ denote that common neighbour of $p_{r_1},p_{r_2}$. Let $X' = \{x'_1,\dots, x'_{\left\lfloor {\frac{|Z_{\delta}|- (k-1)^2}{2}} \right\rfloor}\}$. If $X'$ contains $k$ distinct vertices, then without loss of generality $x'_1 \neq \ldots \neq x'_k$. For every $1 \leq r \leq k$, there is some $C_{2\ell}$ containing $x_0, x'_r$ and vertices from $\{p_{r_1}, \ldots, p_{r_2}\}$, given the pairs $(p_{r_1},p_{r_2})$ have pairwise distance at least $1$ this results in \forbiddenGraph. Therefore, by the pigeon hole principle there is some vertex, we call this $x_{\delta+1}$, which is the common vertex for at least $\left\lfloor {\frac{|Z_{|\delta}|- (k-1)^2}{2(k-1)}} \right\rfloor$ different pairs. Note that $x_{\delta+1}$ is adjacent to at least $\left\lfloor {\frac{|Z_{\delta}|- (k-1)^2}{2(k-1)}} \right\rfloor$ vertices in $\hat{Z}_{\delta}$, and $\left\lfloor {\frac{|Z_{\delta}|- (k-1)^2}{2(k-1)}} \right\rfloor$ vertices in $Z_{\delta}$ that is properties $1$, $3$ and $5$ hold.

        Given $|Z_0| \geq 4^kk^{k+1} \geq 4^{k-1}(k-1)^{k-1}\cdot ((k-1)^2+1)$ there exist vertices $x_0, \ldots, x_{k-1}$, and sets $Z_{0} \supseteq \ldots \supseteq Z_{k-1}$ and $\hat{Z}_{0} \supseteq \ldots \supseteq \hat{Z}_{k-1}$ with each of the properties described above. From property $5$, for every $1 \leq i \leq k-1$, $N(x_{i}) \cap \hat{Z}_{i-1} \neq \emptyset$, this implies there is some vertex in $\hat{Z}_{i-1}$ adjacent to $x_i$, let $y_i$ be some arbitrary vertex in $\hat{Z}_{i-1}$ adjacent to $x_i$. Let $y_k$ be some arbitrary vertex in $\hat{Z}_{k-1}$. From property $4$, $\hat{Z}_i\cap N(x_0,\dots, x_{i})=\emptyset$ meaning, vertices $y_i$ are distinct for each $1 \leq i \leq k$. Further from property $6$, $\{y_1, \ldots, y_k\} \subseteq Z^+_0$. That is for every $1 \leq i \leq k$ there is some distinct $p_j \in Z_0$, we call this vertex $y'_i$, such that $p_{j+2\ell-4} = y_i$. Let $Y_i = \{p_{j}, \ldots, p_{j+2\ell-4}\}$ where $j$ is chosen to satisfy $p_{j+2\ell-4}=y_i$. Note the sets $Y_1, \ldots, Y_k$ are pairwise disjoint.

        Further note $|Z_{k-1}| \geq 2k$, meaning that there exist a set of vertices $C = \{c_1,\ldots,c_k\} \subseteq Z_{k-1}$ such that $y'_i \notin C$ for any $1 \leq i \leq k$.
        Let $x_k$ be the common neighbour of $y_k$ and $c_k$, as $y_k \in \hat{Z}_{k-1}$, $x_k \neq x_i$ for any $0 \leq i \leq k-1$. However, this is a contradiction as for each $1 \leq i \leq k$, there is some cycle of length $2\ell$ containing $x_0$, $c_i$, $x_i$ and $Y_i$, given these cycles have a single common vertex, $x_0$, this describes \forbiddenGraph.
    \end{claimproof}

    We divide $P$ into pairwise disjoint subpaths each containing $m'$ vertices, where $m' = (2\ell-2)((k-1)^2+1)+\ell \geq 10\ell k^2$. We call each of these subpath a segment of $P$. As $m = 4^{k+1}k^{k+4}\ell$, there are at least $4^{k+1}k^{k+2}$ such segments. Let $x_0$ be the common neighbour of $p_0$ and $p_{2(\ell-1)}$. From Claim~\ref{clm:vbndnei}, $x_0$ cannot have neighbours in $4^kk^{k+1}$ different segments (excluding the final segment, i.e. that containing $p_m$). Given there are at least $4^kk^{k+1}+1$ segments there is some segment which does not contain a neighbour of $x_0$.

    \begin{mclaim}\label{clm-seg}
        Let $X$ be a set of external vertices and $Q = (q_0, \cdots, q_{m'})$ be some segment which does not contain a neighbour of any vertex in $X$. Then there is a $C_{2\ell}$ containing $p_0$ and exactly $2$ external vertices $ x\notin X$ and $x' \notin X$ and vertices from $Q$.
    \end{mclaim}
    \begin{claimproof}
        Let $x$ be the common neighbour of $p_0$ and $q_0$. Let $y$ the common neighbour of $p_0$ and $q_{2\ell-4}$. As both $x$ and $y$ have a neighbour in $Q$ we know that $x\notin X$ and $y\notin X$ by assumption. If $y \neq x$, we let $x'=y$, our claim holds as the cycle $(p_0,x,q_0, \dots, q_{2\ell-4},x')$ is of length $2\ell$ and contains $p_0$ and exactly $2$ external vertices $ x\neq x_0$ and $x' \neq x_0$ and vertices from $Q$. Otherwise $y = x$.

        
        Repeating this argument, as $m' > 2(2\ell-4)((k-1)^2+1)+\ell$, either there is some $1 \leq \delta \leq 2((k-1)^2+1)$ and vertex $x' \neq x$, $x' \notin  X$ such that $x$ is adjacent to $q_{\delta(2\ell-4)}$ and $x'$ is the common neighbour of $q_{(\delta+1)(2\ell-4)}$ and $p_0$ or $q_{\delta(2\ell-4)}$ is adjacent to $x$ for all $0 \leq \delta \leq 2((k-1)^2+1)$. In the former case, we have found a cycle $(p_0, x, q_{\delta(2\ell-4)}, \dots, q_{(\delta+1)(2\ell-4)}, x')$ of length $2\ell$ containing $p_0$, two external vertices $x\notin X$, $x'\notin X$ and vertices from $Q$ as claimed.  On the other hand, the latter case contradicts Claim~\ref{clm:dist2l-4} as $x$ has $(k-1)^2+1$ pairs of neighbours with pairwise distance at least $\ell$. 
    \end{claimproof}

    In the following we recursively define $k-1$ distinct segments $Q_1,\dots, Q_{k-1}$ and $2k-2$ distinct external vertices $x_1, \dots, x_{k-1}$ and $x_1',\dots, x_{k-1}'$  such that for every $i\in [k-1]$ we have that $x_i\neq x_0$, $x_i'\neq x_0$ and there is a cycle of length $2\ell$ containing $p_0$, $x_i, x_i'$ and vertices from $Q_i$.

    As highlighted above, there are at least $4^kk^{k+1}+1$ segments, from Claim~\ref{clm:vbndnei}, and hence there is some segment which does not contain a neighbour of $x_0$. Let $Q_1$ be some segment which does not contain any neighbour of $x_0$, from Claim~\ref{clm-seg} there is some pair of vertices, call these $x_1, x'_1 \neq x_0$ such that there is a cycle of length $2\ell$ containing $p_0$, $x_1, x_1'$ and vertices from $Q_1$.
    Suppose now there are segments $Q_1,\dots, Q_{\delta}$ and distinct external vertices $x_1, \dots, x_{\delta}$ and $x_1',\dots, x_{\delta}'$ for some $\delta < k-1$. Let $X_{\delta} = \{x_0,x_1, \dots, x_{\delta}, x_1',\dots, x_{\delta}'\}$. From Claim~\ref{clm:vbndnei}, no vertex in $X_{\delta}$ can have neighbours in $4^kk^{k+1}$ different segments, as $|X_{\delta}| = 2\delta+1$ and there are at least $(2\delta+1) \cdot 4^kk^{k+1}+1$ segments, there is some segment which does not contain a neighbour in $X$, let $Q_{\delta+1}$ be such a segment. Again from Claim~\ref{clm-seg} there is some pair of vertices, call these $x_{\delta+1}, x'_{\delta+1} \notin X_{\delta}$ such that there is a cycle of length $2\ell$ containing $p_0$, $x_{\delta+1}, x'_{\delta+1}$ and vertices from $Q_{\delta+1}$.

    Given segments $Q_1,\dots, Q_{k-1}$ and vertices $x_1, \dots, x_{k-1}$ and $x_1',\dots, x_{k-1}'$ as described above, there are $k-1$ cycles with a common vertex $p_0$. Further the cycle $(x_0, p_0, \ldots, p_{2(\ell-1)})$ has length $2\ell$ and contains no vertex from the segments $Q_1,\dots, Q_{k-1}$, as these segments do not contain a neighbour of $x_0$. That is, this describes \forbiddenGraph. \qed
\end{proof}

\end{document}
